\section{Related Work}
\label{sec:related_work}


Automated data-driven discovery has been a long-standing dream of AI \cite{majumder2024data,kitano2016artificial}.
Although there have been a range of {\bf data-driven discovery systems},
from early ones that fit equations to idealized data, e.g., Bacon \citep{Langley1981DataDrivenDO},
to more modern ones handling complex real-world problems, e.g., AlphaFold \cite{jumper2021highly},
their associated datasets are task-specific and customized to a
pre-built pipeline. In contrast, \discoverybench{} aims to be a general
test over multiple tasks, including testing whether systems can
design appropriate pipelines themselves.

A number of datasets and tools are available for {\bf AutoML}, a related
technology aimed at automating workflows for building optimal machine
learning models \cite{Jin2023AutoKerasAA,Zhang2023AutoMLGPTAM,LeDell2020H2OAS}.
AutoML tools include packages like Scikit \cite{feurer2015efficient},
and embedded in cloud platforms such as Google Cloud Platform,
% \footnote{\url{cloud.google.com/automl}}
Microsoft Azure,
% \footnote{\url{azure.microsoft.com/en-us/products/machine-learning/automatedml/}}
and Amazon Web Services.
% \footnote{\url{aws.amazon.com/machine-learning/automl/}}
However, associated datasets for AutoML are primarily used for training models,
rather than for open-ended discovery tasks.

Similarly, there are several datasets that test {\bf statistical
analysis} in various fields, e.g., \cite{Shao2023ExploringPI,Li2024AutomatedSM,Yang2022ASL}.
Software packages like Tableaux, SAS, and R also support users in that task.
However, these datasets and tools are designed specifically for data analysis,
while \discoverybench{} aims to automate the broader pipeline including ideation,
semantic reasoning, and pipeline design, where statistical analysis is just one component.

% Reasoning and understanding over tables has been a prominent field entailing 
One recent dataset similar in spirit to ours is QRData \cite{liu2024llms}.
QRData also explores LLM capabilities but targets statistical/causal analysis for well-defined (mainly) textbook questions that have unique, (mainly) numeric gold answers.
In contrast, \discoverybench{} has no prescribed boundaries on statistical techniques to apply,
uses open-ended questions and answers, and complex tasks drawn from state-of-the-art published work.


